\scanpage{78}
``nul en codim $<q-1$''. Par l'hyp. de récurrence
il s'ensuit que $H^p(Y,R^qf_*(F))\ne0\Rightarrow$
\margin{$n-q+1$}
$p\leq n-(q-1)$ \unsure{ou} i.e. $p+q\leq n+1$. Donc $E_2^{pq}=0$
sauf si $p+q\leq n+1$, d'où la conclusion\ldots

Par suite, les th. 1 et 2 sont prouvés
dans la situation \deleted{``\unsure{ess.|loc.} t.f. sur}
``\deleted{ponct. \unsure{ess.|est}}\inserted{\unsure{E-ponctuellt|E-ponctuellement}} t.f. sur $k$''. Il en est
donc de \unsure{$=$} pour le lemme 3, de \unsure{$=$}
pour ``ess. de t.f. sur $k$''.
% slanted
\margin{\uline{\unsure{Concl.}}}

Nous allons maintenant le prouver
\unsure{également|a fortiori} quand il s'agit de préschémas
\unsure{E-ponctuellement} de t.f. sur un
anneau \inserted{\uline{noeth}} $\Lambda$ \uline{de dim $\leq1$}. Nous le
prouvons sous \uline{\unsure{sa|ses} forme locale} $2_n$,
par \uline{récurrence sur} $n$, trivial si $n=0$,
on peut supposer $n>0$, et le th. démontré
pour $n'<n$. On peut supposer $\Lambda$ \uline{str. local}
\inserted{intègre} \uline{de dim 1}, $\Lambda\to A$ \deleted{et} \inserted{local} \deleted{intègre \uline{dominant} $\Lambda$,} \inserted{injectif.}
% slanted
\margin{\uline{Arranger la} \uline{démonstration}}
% slanted
\margin{\uline{Récurrence} \uline{inutile}}
% archive sheet number: 77
